\section{Discussion}
\label{sec:discussion}

\subsection{Information Sufficiency, Not Token Budgets}

The convergence of three independent experiments (the core saturation study,
the randomised ablation, and the padding control) supports a single
mechanism: quality saturates once \emph{task-critical information} has been
communicated.
The saturation \emph{shape} (quality rising then plateauing) is robust across
prompt orderings, but the specific token count at which it occurs is not: it
shifts with which layers appear early.
Irrelevant tokens at matched counts produce flat or declining quality.
The practical implication is that prompt optimisation should focus on
\emph{which mechanisms to include}, not on \emph{how many tokens to budget}.

\subsection{Schema Gating: One Layer Unlocks Structured Tasks}

The strongest result of the study is that, on deterministic ground truth,
structured tasks saturate at a single \emph{schema-resolving} layer, the layer
that tells the model what shape the answer must take.
Four lines of evidence converge:

\begin{enumerate}
  \item \textbf{Concentration on ground truth.} One layer captures the bulk of
    all quality gain: classification 87\% (task label, L2), product extraction
    92\% (JSON format, L3), NER 95\% (format, L3). Later layers (definitions,
    persona, guidelines, and even the worked example) add little.
  \item \textbf{Out-of-sample prediction.} We predicted both the class
    (schema-gated) and the exact layer (L3 format) for the held-out NER task
    \emph{before} running it; ground truth confirmed both. The mechanism has
    predictive power, not just descriptive fit.
  \item \textbf{Step shape.} Product extraction and NER are near-perfect step
    functions ($q \approx 0$ until the format layer, then $q > 0.9$). The
    model is not gradually improving; it is switching from ``cannot produce
    parseable output'' to ``can.''
  \item \textbf{Knowledge-gated contrast.} QA and math reach ceiling at L1
    precisely because they require no schema: the answer is a fact or a number
    the model already holds, and they stay flat even on hard examples.
\end{enumerate}

This sharpens the earlier ``information sufficiency'' framing into a concrete,
testable mechanism: for structured tasks the task-critical information is the
\emph{output schema}, and once it is communicated, elaboration stops helping.
We note the judge would have obscured this: it credited the worked example and
implied a ``distributed'' profile, which ground truth overturns
(Section~\ref{sec:result-judge}).
A direct mechanistic-interpretability test of how models internalise the format
layer remains future work.

\subsection{Harmful Prompt Content}

The math L3 finding demonstrates that prompt content can be actively harmful:
instructing models to ``give only the final numerical answer'' suppresses
chain-of-thought and lowers quality ($-0.366$ on the LLM judges, $-0.071$ on
exact-match, across all seven models).
The metric gap is itself a finding: the judge penalises missing reasoning
heavily, while exact-match registers only the cases where suppression also
changed the answer; both agree on the sign.
This extends beyond the ``diminishing returns'' framing: certain instructions
are not merely unhelpful but counterproductive.
The standard sigmoid saturation model does not capture tasks where
intermediate prompt layers are detrimental, and this non-monotonicity is
important for practitioners to be aware of.

\subsection{Practical Guidelines}

The three-class account yields concrete, task-specific guidance:

\begin{itemize}
  \item \textbf{Schema-gated (classification, extraction, NER):} Spend the
    prompt budget on \emph{one} well-formed output-schema specification (the
    label space or a minimal format template) and stop. This single layer
    captures 87--95\% of the achievable gain; worked examples and personas add
    little once the format is clear ($\approx$40--100 tokens suffice). The
    cheapest high-value prompt is a format spec, not a worked example.
  \item \textbf{Difficulty-gated (instruction following):} Elaboration pays off
    only in proportion to task difficulty. For simple single-constraint tasks,
    minimal prompts suffice; for multi-constraint tasks, add the definition and
    guideline layers, which become load-bearing as constraints accumulate.
  \item \textbf{Knowledge-gated (QA, math, summarisation):} The bare question
    suffices. Invest in retrieval (for QA) or chain-of-thought formatting (for
    math) rather than instruction elaboration, and \emph{never} suppress
    step-by-step reasoning, which costs $-0.366$ quality on math.
\end{itemize}

\subsection{Relationship to Prompt Compression}

Our work and prompt compression methods (LLMLingua, LLMLingua-2) operate at
different stages of the prompt lifecycle.
Compression decides what to keep \emph{given a written prompt}; our marginal
contribution analysis decides what to \emph{write in the first place}.
For classification, LLMLingua-2 might compress L7 ($\approx$180 tokens) to
$\approx$36 tokens by removing filler, but our analysis shows the task label
alone (L1$\to$L2, $\approx$5 tokens) already captures 87\% of the quality gain
on ground truth, so the L7 elaboration was largely \emph{unnecessary content}
rather than redundantly worded content.
Compression cannot recover this: it shortens text that should never have been
written.
The two approaches are complementary and stack: author the right mechanisms
(our work), then compress the wording within them (prompt compression).

\subsection{Limitations}
\label{sec:limitations}

We foreground the limitations that genuinely qualify our claims before the
routine ones.

\paragraph{The ``insensitive'' tasks are only partly resolved.}
QA is flat under both metrics and stays flat on hard HotpotQA examples, so we
are confident it is knowledge-gated. Math is murkier: it is knowledge-gated for
strong models but edges toward difficulty-gated for weaker ones on hard GSM8K
problems, so ``insensitive'' is a model-dependent statement, not a task
property. The below-ceiling stratification that first suggested this is
confounded by regression to the mean and should not be read as independent
evidence; only the hard-example experiment is.

\paragraph{Curve fitting under-detects concentrated tasks.}
With $n = 7$ points per curve, the F-test detects gradual sigmoids well but
under-detects step functions: classification is significant for only 2/7 models
despite an 87\% L2 jump. We therefore do not rely on per-curve significance for
schema-gated tasks; the marginal-contribution analysis and the 7/7-significant
product-extraction fits (which \emph{are} smooth enough to fit) carry the
claim. Finer level spacing (15--20 levels) would let curve fitting capture
step-shaped tasks too.

\paragraph{Classification domain breadth.}
Classification draws from sentiment only. Product extraction and the held-out
NER task provide two further schema-gated tasks with different output
structures, but extending to additional classification domains (topic, intent,
specialised label spaces such as ICD-10) would test whether the concentrated
shape persists when the label space is harder to infer.

\paragraph{Single-turn scope.}
The study focuses on single-turn instruction tasks.
Agentic, RAG, and multi-turn settings may exhibit different saturation
dynamics, and our mechanism-level findings yield testable predictions for
these settings (e.g., tool-call format specifications may exhibit
concentrated sensitivity analogous to classification).

\paragraph{English only.}
All prompts and examples are in English.
Saturation behaviour may differ for low-resource languages where prompt
elaboration is more necessary to disambiguate tasks.

\paragraph{Model deprecation and versioning.}
Two models were affected by provider deprecations across the study, a
reproducibility constraint rather than a design flaw.
kimi-k2 was deprecated by Groq between the main study and the later
experiments, producing 404 errors on all 490 ablation attempts; its
immutable main-study outputs are retained, and the ablation and other
follow-on analyses are restricted to the models with complete data.
gemini-2.0-flash was likewise deprecated, so the later experiments
(held-out NER, hard examples, sampling stability) use gemini-2.5-flash in its
place.
The version change shifts effect \emph{magnitudes} (classification L1$\to$L7
gain rises from $+0.110$ on gemini-2.0-flash to $+0.345$ on gemini-2.5-flash),
but the \emph{direction} of every effect is preserved across both versions, so
no qualitative conclusion depends on a single model build.

\paragraph{Taxonomy validated on one held-out task, not many.}
The three-class account (schema-, difficulty-, knowledge-gated) was abstracted
from six tasks, but the schema-gated class now has out-of-sample support: we
predicted NER's class and gain locus before running and were confirmed. This is
one held-out task, not a battery: the difficulty- and knowledge-gated classes
have not been validated out of sample, and a single confirmation cannot rule
out a task that defies the scheme. Broader held-out validation (code generation,
data-to-text, dialogue-state tracking) is the natural next step.
